\section{Introduction}

Late-time cosmological tensions and directional residuals motivate tests that are capable of separating local structure, survey geometry, and genuinely new large-scale dynamics. The scalar-topographic program introduced a compact hyperspherical parent \(S^3(R_H)\) and a long-wavelength scalar response whose largest observational effect occurs at low redshift \citep{Carberry2026Topographic}. A later covariant formulation placed the topographic scalar in the cubic-Galileon sector of Horndeski theory, preserving second-order time evolution while allowing nontrivial late-time scalar--metric braiding.

The present work asks a narrower question: can the long-wavelength topographic response be represented as the large-scale branch of the same constraint-reduced response logic used in the Berger--Hopf program, without assuming that the full microscopic BHSM construction has already been compactified into cosmology? We answer this at the level of a controlled bridge model.

The central mathematical chain is
\[
S[\Phi]
\longrightarrow
\widehat{\mathcal H}_{\rm phys}
\longrightarrow
\mathcal K_{\rm topo}
\longrightarrow
\lambda_{\rm topo}
\longrightarrow
T_n
\longrightarrow
\delta g_{\mu\nu}
\longrightarrow
\{D_L,H,f\sigma_8,\Phi+\Psi\}.
\]

The distinction between established input and bridge extension is essential. The \(S^3\) harmonic spectrum and cubic-Galileon background are treated as the cosmological starting point. The Schur/Feshbach reduction is an exact mathematical operation once a physical complement is specified. The identification of its low-energy coefficients with the topographic response is a bridge assumption to be tested, not a declaration that the complete microscopic BHSM action has already been reduced to four-dimensional cosmology.

The paper makes four main contributions. First, it standardizes the \(S^3\) harmonic convention and shows how the observed sky dipole arises from the four-dimensional \(n=1\) eigenspace rather than by identifying a radial mode with an angular dipole. Second, it derives a reduced kernel in which a positive higher-spatial-gradient term can emerge after integrating out a stable complement. Third, it identifies a parameter interval in which the \(n=1\) response is selectively softened while all \(n\ge2\) harmonics are stiffened. Fourth, it gives an executable preregistration protocol for curvature, supernova anisotropy, growth, and lensing tests.

The claim boundary is therefore:
\begin{quote}
This work does not assume that the scalar-topographic cosmological field has already been derived from the complete microscopic Berger--Hopf action. Exact reductions are distinguished from cross-scale identifications and exploratory reference-branch assumptions. Prospective observables are frozen before comparison with designated independent data.
\end{quote}
